% THIS DOCUMENT WAS COMPUTER GENERATED FROM  diploma-prazno
% IT REQUIRES TO BE RUN FROM TERMINAL SINCE GLOBAL .vsc SETTINGS

\documentclass{beamer}

\usetheme{Warsaw}

\usepackage[slovene]{babel}
\usepackage[utf8]{inputenc}
\usepackage[T1]{fontenc}
\usepackage{lmodern}

\usepackage{amsfonts, amsmath, amsthm}
\usepackage{enumitem}

% --- Teorematična okolja ---
\newtheorem{izrek}{Izrek}
\newtheorem{trditev}[izrek]{Trditev}
\newtheorem{posledica}[izrek]{Posledica}
\newtheorem{definicija}[izrek]{Definicija}

% --- Ukazi ---
\newcommand{\red}{\operatorname{red}}
\newcommand{\srg}{\operatorname{srg}}
\newcommand{\diam}{\operatorname{diam}}
\newcommand{\N}{\mathbb{N}}
\newcommand{\ZZ}{\mathbb{Z}}
\newcommand{\reach}{\operatorname{Reach}}
\newcommand{\gr}{\mathscr{G}}

\newcommand{\A}{\mathcal{A}}
\newcommand{\II}{\mathbf{I}}
\newcommand{\IIw}[2]{#1 \;\II\; #2}
\newcommand{\IIwd}[2]{#1 \;\II^D\; #2}
\newcommand{\PP}{\mathcal{P}}
\newcommand{\LL}{\mathcal{L}}

% --- Naslov ---
\title{Moorejevi grafi in posplošeni večkotniki}
\author{Luka Lavš}
\date{\today}

\begin{document}

% --- Title ---
\begin{frame}
    \titlepage
\end{frame}

% ------------------------------
\section{Moorejevi grafi}

\begin{frame}
\frametitle{Definicija Moorejevega grafa}
\begin{definicija}
Naj bo $G=(V,E)$ graf s premerom $D$ in maksimalno stopnjo $\Delta$.  
Če velja
\[
|V| = 1 + \Delta + \Delta(\Delta-1) + \dots + \Delta(\Delta-1)^{D-1},
\]
pravimo, da je $G$ \emph{Moorejev graf}.
\end{definicija}
\end{frame}

\begin{frame}
\frametitle{Lastnosti}
\begin{trditev}
Moorejev graf $G$ je $\Delta$-regularen. Med vsakima dvema vozliščema obstaja natanko ena najkrajša pot dolžine največ $D$.
\end{trditev}

\begin{posledica}
Ožina Moorejevega grafa je $g = 2D + 1$.
\end{posledica}
\end{frame}

\begin{frame}
\frametitle{Primeri Moorejevih grafov}
\begin{itemize}
    \item $D=2, \Delta=2$: pentagon
    \item $D=2, \Delta=3$: Petersonov graf
    \item $D=2, \Delta=7$: Hoffman-Singletonov graf
    \item $D=2, \Delta=57$: odprto vprašanje
\end{itemize}
\end{frame}

\begin{frame}
\frametitle{Matrična lastnost}
Za regularen Moorejev graf $G$ stopnje $\Delta$ z matriko sosednosti $A$ velja
\[
A^2 + A - (\Delta-1)I = J,
\]
kjer je $I$ identiteta in $J$ matrika samih enic.
\end{frame}

\begin{frame}
\frametitle{Spodnja meja za graf z ožino $g$}
\begin{trditev}
Naj bo $G$ $\Delta$-regularen z ožino $g$. Tedaj velja
\[
|V(G)| \ge
\begin{cases}
\frac{\Delta(\Delta-1)^{(g-1)/2}-1}{\Delta-2}, & g \text{ liho},\\[1em]
2\frac{(\Delta-1)^{g/2}-1}{\Delta-2}, & g \text{ sodo}.
\end{cases}
\]
\end{trditev}
\end{frame}

% ------------------------------
\section{Geometrije ranga 2}

\begin{frame}
\frametitle{Definicija geometrije}
Geometrija \(\Gamma=(\A_1, \dots, \A_n, \II)\) je nabor nepraznih, paroma disjunktnih množic \(\A_i\) in incidenčne relacije
\[
\II \subseteq \bigcup_{i\neq j} \A_i \times \A_j.
\]
\end{frame}

\begin{frame}
\frametitle{Geometrija ranga 2}
\begin{definicija}
Geometrija ranga 2 je \(\Gamma=(\PP, \LL, \II)\), kjer
\begin{itemize}
    \item \(\PP\) množica točk
    \item \(\LL\) množica črt
    \item $l \in \LL$ je v relaciji z $p \in \PP$ natanko tedaj, ko $l$ in $p$ incidirata
\end{itemize}
\end{definicija}
\end{frame}

\begin{frame}
\frametitle{Incidenčni graf}
\begin{definicija}
Naj bo \(\Gamma=(\PP,\LL,\II)\).  
Incidenčni graf \(G\) geometrije \(\Gamma\) ima
\[
V(G) = \PP \cup \LL, \quad E(G) = \{ uv \mid u \text{ in } v \text{ incidirata} \}.
\]
\end{definicija}
\end{frame}

% ------------------------------
\section{Posplošeni večkotniki}

\begin{frame}
\frametitle{Večkotnik dolžine $n$}
\begin{definicija}
Naj bo \(\Gamma=(\PP,\LL,\II)\). Večkotnik dolžine $n$ je zaporedje \(x_1, \dots, x_{2n} \in \PP \cup \LL\) z lastnostmi:
\begin{enumerate}[label=(\roman*)]
\item elementi se izmenjujejo po vrstah (\(x_{2i-1} \in \PP\), \(x_{2i} \in \LL\))
\item vsak element incidirata s sosednjim elementom (\(x_i\) incidirata z \(x_{i+1}\)), pri čemer \(x_{2n+1}=x_1\)
\end{enumerate}
\end{definicija}
\end{frame}

\begin{frame}
\frametitle{Posplošeni večkotniki in Moorejevi grafi}
\begin{definicija}
Posplošeni večkotnik je struktura, ki generalizira koncept večkotnika in je ključna za povezavo Moorejevih grafov lihe ožine z incidenčnimi grafi.
\end{definicija}

\begin{itemize}
\item Primer: incidenčni graf PG(2,2) ustreza Petersonovemu grafu
\item Hoffman-Singletonov graf izhaja iz posplošenega petkotnika
\item Moorejevi grafi z liho ožino ustrezajo incidenčnim grafom posplošenih večkotnikov
\end{itemize}
\end{frame}

\begin{frame}
\frametitle{Odprta vprašanja}
Obstoj Moorejevega grafa valence \(\Delta=57\) še ni dokazan.  
Pogoste metode konstrukcije (greedy approach) pri \(\Delta \ge 7\) ne uspejo.
\end{frame}

\end{document}
